% =============================================================================
% MLSys·im Tutorial Tutorial — Parts 0–4 (Morning Session)
% =============================================================================
\documentclass[aspectratio=169, 12pt]{beamer}
\usepackage{../../../slides/assets/beamerthememlsys}

\mlsyssetup{
  volume       = {Tutorial},
  chapter      = {Module 2},
  logo         = {../../../slides/assets/img/logo-mlsysbook.png},
  instlogo     = {../../../slides/assets/img/logo-harvard.png},
  chaptertitle = {MLSys·im: Advanced Single-Node Analysis},
}

% --- Fonts ---
\usepackage[T1]{fontenc}
\usepackage[scaled=0.9]{helvet}
\usepackage{courier}
\renewcommand{\familydefault}{\sfdefault}

% --- Packages ---
\usepackage{amsmath}
\usepackage{booktabs}
\usepackage[table]{xcolor}
\usepackage{listings}
\usepackage{tikz}
\usetikzlibrary{arrows.meta, positioning, calc, decorations.pathreplacing}

% --- Code listings ---
\lstset{
  language=Python,
  basicstyle=\ttfamily\footnotesize,
  keywordstyle=\color{crimson}\bfseries,
  stringstyle=\color{datastroke},
  commentstyle=\color{midgray}\itshape,
  backgroundcolor=\color{computeblue!20},
  frame=single,
  rulecolor=\color{computestroke},
  numbers=none,
  breaklines=true,
  columns=fullflexible,
  keepspaces=true,
  showstringspaces=false,
  xleftmargin=4pt,
  xrightmargin=4pt,
  aboveskip=6pt,
  belowskip=4pt,
}

% --- Convenience macros ---
\newcommand{\mlsysim}{\texttt{mlsysim}}
\newcommand{\wallbox}[2]{%
  \begin{block}{#1}#2\end{block}%
}
\newcommand{\PredictStart}{\begin{alertblock}{Predict Before You Peek}}
\newcommand{\PredictEnd}{\end{alertblock}}

% --- Image paths ---
\graphicspath{{images/}}

% --- Section count (must match actual \section{} count) ---
\setcounter{mlsystotalsections}{6}

\title{MLSys·im Tutorial --- Module 2}
\subtitle{Advanced Single-Node Analysis}
\author{Vijay Janapa Reddi}
\institute{Harvard University}
\date{Tutorial}

% =============================================================================
\begin{document}

% --- Roadmap ---
\begin{frame}{Roadmap: Conference Tutorial}
\centering\small
\begin{tabular}{rll}
\toprule
\textbf{Module} & \textbf{Topic} & \textbf{Status} \\
\midrule
Module 1 & Foundations \& Architecture & \checkmark Done \\
\rowcolor{crimson!12}
\textbf{Module 2} & \textbf{Advanced Single-Node Analysis} & \textbf{$\leftarrow$ You are here} \\
Module 3 & Scale, Dollars, and Carbon & \\
Module 4 & Design Space Exploration \& Synthesis & \\
\bottomrule
\end{tabular}
\end{frame}


\section{Memory Walls \& Serving}

\begin{frame}{Beyond FLOPs: The Data Wall}
GPUs are so fast they often starve waiting for data. This occurs in two places:
\begin{itemize}
  \item \textbf{Ingestion (I/O):} Can the storage sub-system push bytes fast enough?
  \item \textbf{Transformation (CPU):} Can the CPUs decode JPEGs and augment data fast enough?
\end{itemize}
\vfill
\begin{block}{The Algebra of Data Stalls}
\[
\text{Ingestion Stalled} \iff \text{Demand BW} > \text{Storage BW}
\]
\[
\text{Transformation Stalled} \iff \frac{\text{Batch Size}}{\text{CPU Throughput}} > T_{\text{GPU\_Step}}
\]
\end{block}
\end{frame}

\begin{frame}[fragile]{Live Demo: Uncovering the Data Wall (Scenario C)}
\begin{lstlisting}[language=Python]
from mlsysim.solvers import DataModel, TransformationModel
from mlsysim.hardware.registry import Hardware
from mlsysim.core.constants import Q_

demand_rate = Q_("40000 1/s") * Q_("150 KB") # ~6 GB/s

data_result = DataModel().solve(
    workload_data_rate=demand_rate, hardware=Hardware.Cloud.A100)
print(f"I/O Stalled: {data_result.is_stalled}") # False

transform_result = TransformationModel().solve(
    batch_size=40000,
    cpu_throughput_per_worker_hz=850, # JPEG decode + crop
    num_workers=64                    # 8 CPUs per GPU
)
print(f"CPU Stalled: {transform_result.is_bottleneck}") # True
\end{lstlisting}
\end{frame}


% --- Slide 2.1: Key Question ---
\begin{frame}{Key Question}
\note{[1 min] ``You have seen that LLM inference at batch 1 is
memory-bound. But serving is more complex than a single forward
pass. What makes LLM serving fundamentally different from
CNN inference?''
% -- FLEX: [CORE]
}

\centering
\Large\bfseries
Why does the first token take 50\,ms\\[0.2cm]
but each next token only takes 5\,ms?
\end{frame}

% --- Slide 2.2: Two Phases of LLM Serving ---
\begin{frame}{Wall 4: Prefill vs Decode --- Two Different Physics}
\note{[3 min] ``Prefill is like reading a book fast (compute-intensive).
Decode is like looking up one word at a time in a dictionary
(memory-intensive). Same model, different bottlenecks.''
% -- FLEX: [CORE]
}

\small
\wallbox{The Serving Wall}{
\begin{tabular}{lcl}
  \textbf{TTFT} (Prefill) & $=$ & $\dfrac{\text{Prefill FLOPs}}{\text{Peak FLOPS} \times \text{MFU}}$
  \quad\colorbox{computeblue}{Compute-bound} \\[10pt]
  \textbf{ITL} (Decode)   & $=$ & $\dfrac{\text{Weight Bytes}}{\text{Bandwidth}}$
  \quad\colorbox{datagreen}{Memory-bound}
\end{tabular}
}

\vspace{0.3cm}
\begin{columns}[T]
\begin{column}{0.48\textwidth}
\textbf{Prefill} (process the prompt)
\begin{itemize}\setlength\itemsep{2pt}
  \item All prompt tokens in parallel
  \item $O(S^2)$ attention + $O(S \cdot P)$ linear
  \item Compute-bound (high AI)
  \item Determines \textbf{TTFT}
\end{itemize}
\end{column}
\begin{column}{0.48\textwidth}
\textbf{Decode} (generate tokens)
\begin{itemize}\setlength\itemsep{2pt}
  \item One token at a time
  \item Must reload all weights per token
  \item Memory-bound (AI $\approx$ 1)
  \item Determines \textbf{ITL}
\end{itemize}
\end{column}
\end{columns}
\end{frame}

% --- Slide 2.3: KV Cache --- The Hidden Consumer ---
\begin{frame}{Wall 5: The KV Cache --- Hidden Memory Consumer}
\note{[3 min] ``Each active request carries its own memory of the
conversation.'' Quick math: how much KV cache does one Llama-3 8B
request at 4K context need in FP16?
Expected: 2 * 32 * 32 * 128 * 4096 * 2 bytes = ~2 GB.
% -- FLEX: [CORE]
}

\small
\wallbox{The Batching Wall}{
\[
  \text{KV cache} = 2 \times L \times H \times d \times S \times B \times \text{bpp}
\]
}

\vspace{0.2cm}
\textbf{Llama-3 8B at 4K context, FP16:}
\begin{itemize}
  \item $2 \times 32 \times 32 \times 128 \times 4096 \times 2\;\text{bytes} \approx 2\;\text{GB}$ per request
  \item H100 has 80 GB HBM --- model weights take 16 GB
  \item Remaining 64 GB $\div$ 2 GB/request $=$ \textbf{$\sim$32 concurrent requests}
\end{itemize}

\vspace{0.2cm}
\begin{alertblock}{The serving paradox}
  You want high batch size (for throughput) but KV cache limits how
  many requests fit in memory. \alert{Memory capacity, not compute, limits concurrency.}
\end{alertblock}
\end{frame}

% --- Slide 2.4: Live Demo --- ServingModel ---
\begin{frame}[fragile]{Live Demo: Two-Phase Serving Analysis}
\note{[2 min] Run ServingModel live. Point out TTFT vs ITL in output.
Show that TTFT is compute-bound and ITL is memory-bound.
% -- FLEX: [CORE]
}

\small
\begin{lstlisting}
serving = mlsysim.ServingModel()
result  = serving.solve(llama, hw,
                        seq_len=4096, batch_size=1)
print(f"TTFT: {result.ttft:~P}")
print(f"ITL:  {result.itl:~P}")
print(f"KV cache/req: {result.kv_cache_per_request:~P}")
\end{lstlisting}

\vspace{0.2cm}
\begin{exampleblock}{Expected output}
  TTFT $\approx$ 20--50 ms (compute-bound),
  ITL $\approx$ 5 ms (memory-bound),
  KV cache per request $\approx$ 2 GB
\end{exampleblock}
\end{frame}

% --- Slide 2.5: Continuous Batching ---
\begin{frame}{Continuous Batching: Don't Wait, Serve}
\note{[2 min] ``In static batching, the GPU waits for the longest
request to finish. In continuous batching, new requests start
as soon as any slot frees up. Throughput can improve 2--5x.''
% -- FLEX: [OPTIONAL] Can summarize quickly if behind schedule.
}

\small
\begin{columns}[T]
\begin{column}{0.48\textwidth}
\textbf{Static batching}
\begin{itemize}\setlength\itemsep{2pt}
  \item Pad all sequences to max length
  \item GPU idle while short requests finish
  \item Throughput limited by longest request
  \item Simple to implement
\end{itemize}
\end{column}
\begin{column}{0.48\textwidth}
\textbf{Continuous batching}
\begin{itemize}\setlength\itemsep{2pt}
  \item Insert new requests per iteration
  \item No padding waste
  \item Throughput 2--5$\times$ higher
  \item Used by vLLM, TGI, TensorRT-LLM
\end{itemize}
\end{column}
\end{columns}

\vspace{0.3cm}
\centering
\begin{tikzpicture}[scale=0.65, >=Stealth]
  % Static
  \node[font=\scriptsize\bfseries] at (0, 2.5) {Static};
  \foreach \i/\len in {0/4, 1/2, 2/3} {
    \fill[computeblue, draw=computestroke] (0.5, \i*0.7) rectangle ({0.5 + \len*0.5}, \i*0.7+0.5);
    \fill[errorfill, draw=errorstroke, opacity=0.5] ({0.5 + \len*0.5}, \i*0.7) rectangle (2.5, \i*0.7+0.5);
  }
  \node[font=\tiny, errorstroke] at (2.8, 0.7) {waste};

  % Continuous
  \node[font=\scriptsize\bfseries] at (5, 2.5) {Continuous};
  \fill[computeblue, draw=computestroke] (5.5, 1.4) rectangle (7.5, 1.9);
  \fill[datagreen, draw=datastroke] (5.5, 0.7) rectangle (6.5, 1.2);
  \fill[routingorange, draw=routingstroke] (6.7, 0.7) rectangle (7.5, 1.2);
  \fill[computeblue, draw=computestroke] (5.5, 0) rectangle (7, 0.5);
  \fill[datagreen, draw=datastroke] (7.2, 0) rectangle (7.5, 0.5);
  \node[font=\tiny, datastroke] at (8, 0.7) {no waste};
\end{tikzpicture}
\end{frame}

% --- Slide 2.6: PagedAttention ---
\begin{frame}[fragile]{PagedAttention: Virtual Memory for KV Cache}
\note{[2 min] ``Just like OS virtual memory pages physical RAM,
PagedAttention pages the KV cache. Non-contiguous blocks mean
no fragmentation, so you can fit more requests.''
% -- FLEX: [OPTIONAL]
}

\small
\textbf{The problem:} Pre-allocated KV cache wastes memory on short sequences.

\textbf{The solution} (Kwon et al., 2023 --- vLLM):
\begin{itemize}
  \item Divide KV cache into fixed-size \textbf{pages} (e.g., 16 tokens each)
  \item Allocate pages on demand, not up front
  \item Non-contiguous storage eliminates fragmentation
  \item Memory utilization improves from $\sim$50\% to $>$95\%
\end{itemize}

\vspace{0.2cm}
\begin{exampleblock}{Impact}
  With the same 80 GB H100, PagedAttention can serve
  \textbf{2--4$\times$ more concurrent requests} than static allocation.
\end{exampleblock}

\vspace{0.2cm}
\begin{lstlisting}
# mlsysim models this with ContinuousBatchingModel
from mlsysim.core.solver import ContinuousBatchingModel
cb = ContinuousBatchingModel()
result = cb.solve(model, hw, seq_len=4096,
                  max_batch_size=32, page_size=16)
print(f"Max concurrent: {result.max_concurrent_requests}")
\end{lstlisting}
\end{frame}

% --- Slide 2.7: Predict Before You Peek #3 ---
\begin{frame}{Predict: How Many Concurrent Requests?}
\note{[2 min] Predict-before-reveal. Give 60 seconds.
H100 80 GB, Llama-3 8B FP16, 4K context. How many concurrent requests?
Expected: weights = 16 GB, remaining = 64 GB.
PagedAttention (95\% util): ~30. Without (50\% util): ~16.
Turn to your neighbor: did you get the same answer? Why or why not? 60 seconds.
% -- FLEX: [CORE]
}

\centering
\Large\bfseries
H100 (80 GB), Llama-3 8B (FP16), 4K context.\\[0.5cm]
\normalsize
How many concurrent requests can you serve?\\[0.3cm]

\pause

\small
\begin{tabular}{lcc}
\toprule
  & \textbf{Static alloc.} & \textbf{PagedAttention} \\
\midrule
KV cache utilization & $\sim$50\% & $\sim$95\% \\
Effective memory/req & $\sim$4 GB & $\sim$2.1 GB \\
Max concurrent & $\sim$16 & $\sim$30 \\
\bottomrule
\end{tabular}

\vspace{0.3cm}
\pause
\alert{PagedAttention nearly doubles serving capacity without changing hardware.}
\end{frame}

% --- Slide 2.8: Speculative Decoding ---
\begin{frame}[fragile]{Speculative Decoding: Betting on the Draft}
\note{[2 min] ``Use a small draft model to guess the next K tokens.
Then verify all K in a single forward pass of the big model.
If the draft is right 70\% of the time and K=5, you effectively
decode ~3.5 tokens per forward pass instead of 1.''
% -- FLEX: [OPTIONAL]
}

\small
\textbf{Insight:} Decode is memory-bound $\Rightarrow$ the GPU has spare compute.

\textbf{Speculative decoding} (Leviathan et al., 2023):
\begin{enumerate}\setlength\itemsep{2pt}
  \item \textbf{Draft} $K$ tokens with a small model (fast, low quality)
  \item \textbf{Verify} all $K$ tokens in one forward pass of the big model
  \item \textbf{Accept} the longest prefix that matches
  \item Speedup $\approx K \times \alpha$ where $\alpha$ is acceptance rate
\end{enumerate}

\vspace{0.2cm}
\begin{lstlisting}
# mlsysim supports speculative decoding
draft = mlsysim.Models.Language.Llama3_8B  # smaller
result = serving.solve(model, hw, seq_len=4096,
          draft_model=draft, draft_acceptance_rate=0.7)
print(f"Speedup: {result.speculative_speedup:.2f}x")
\end{lstlisting}
\end{frame}

% --- Slide 2.9: Disaggregated Serving ---
\begin{frame}[fragile]{Disaggregated Serving: Right Hardware for Each Phase}
\note{[2 min] ``Split prefill and decode onto different node types.
Prefill nodes optimize for FLOPS, decode nodes optimize for
bandwidth. Transfer KV cache over the network between them.''
% -- FLEX: [OPTIONAL]
}

\small
\textbf{Key insight:} Prefill and decode have \textit{opposite} hardware preferences.

\begin{columns}[T]
\begin{column}{0.48\textwidth}
\textbf{Prefill node}
\begin{itemize}\setlength\itemsep{2pt}
  \item Needs high FLOPS
  \item Moderate memory
  \item E.g., H100 at high utilization
\end{itemize}
\end{column}
\begin{column}{0.48\textwidth}
\textbf{Decode node}
\begin{itemize}\setlength\itemsep{2pt}
  \item Needs high BW
  \item Large memory (for KV cache)
  \item E.g., many smaller accelerators
\end{itemize}
\end{column}
\end{columns}

\vspace{0.3cm}
\begin{lstlisting}
# Disaggregated serving in mlsysim
result = serving.solve(model, hw, seq_len=4096,
          decode_hardware=mlsysim.Hardware.Cloud.A100)
print(f"TTFT: {result.ttft:~P}  ITL: {result.itl:~P}")
\end{lstlisting}

\vspace{0.2cm}
\alert{Trade-off: network transfer of KV cache adds latency between phases.}
\end{frame}

% --- Slide 2.10: Exercise 2 ---
\begin{frame}[fragile]{Exercise 2: Serving Capacity Planning}
\note{[5 min] Attendees work in pairs.
How many concurrent Llama-3 8B requests (4K context) can an H100
serve while maintaining ITL < 10ms?
Expected: ~30 with PagedAttention at FP16. Binding: memory capacity.
Turn to your neighbor: did you get the same answer? Why or why not? 60 seconds.
% -- FLEX: [CORE]
}

\centering
\Large\bfseries
Hands-On Exercise\\[0.5cm]
\normalsize

\textbf{Question:} You run Llama-3 8B (FP16) on one H100 with 4K context.\\
Your SLA requires ITL $<$ 10 ms.\\
How many concurrent requests can you serve?

\vspace{0.3cm}
\begin{lstlisting}
from mlsysim.core.solver import ContinuousBatchingModel
cb = ContinuousBatchingModel()
result = cb.solve(llama, hw, seq_len=4096,
                  max_batch_size=64, page_size=16)
print(f"Max concurrent: {result.max_concurrent_requests}")
\end{lstlisting}

\vspace{0.2cm}
\textit{Bonus: What happens if you switch to INT8 precision?
Does concurrency double?}
\end{frame}

% --- Slide 2.11: Fallacies ---
\begin{frame}{Fallacies: Serving Edition}
\note{[2 min] Walk through each fallacy with the quantitative
counter-evidence.
% -- FLEX: [OPTIONAL] --- can trim to 2 fallacies if behind.
}

\small
\textbf{Fallacy:} \textit{``Faster GPUs always reduce latency.''}\\
A 3.2$\times$ FLOPS improvement (A100 $\to$ H100) yields only
1.7$\times$ ITL improvement because decode is memory-bound.

\vspace{0.3cm}
\textbf{Fallacy:} \textit{``Doubling memory doubles serving capacity.''}\\
Weights are fixed overhead. Going from 80 GB to 160 GB adds 80 GB,
but KV cache per request stays $\sim$2 GB. Capacity goes from $\sim$32
to $\sim$72 ($\sim$2.3$\times$), not 2$\times$.

\vspace{0.3cm}
\textbf{Fallacy:} \textit{``Batch size 1 is fine for LLM inference.''}\\
At bs=1, MFU $<$ 5\%. You are paying for 989 TFLOPS and using $<$ 50.
Continuous batching can recover 10--20$\times$ throughput.
\end{frame}

% --- Slide 2.12: Part 2 Key Takeaway ---
\begin{frame}{Part 2: Key Takeaway}
\note{[1 min] One sentence summary, repeat twice.
% -- FLEX: [CORE]
}

\centering\Large

\textbf{LLM serving has two phases\\with opposite bottlenecks.}\\[0.5cm]

\normalsize
\begin{itemize}
  \item \textbf{Prefill} (TTFT) is compute-bound --- optimize with parallelism.
  \item \textbf{Decode} (ITL) is memory-bound --- optimize with bandwidth.
  \item \textbf{KV cache} limits concurrency --- optimize with PagedAttention.
  \item \texttt{ServingModel.solve()} decomposes both phases in one call.
\end{itemize}
\end{frame}

% =============================================================================
% PART 3: COMPRESSION & EFFICIENCY (8 slides)
% =============================================================================

\section{Compression \& Efficiency}

% --- Slide 3.1: Key Question ---
\begin{frame}{Key Question}
\note{[1 min] ``If we can shrink the model, we move less data,
and the memory wall recedes. But there is a catch.''
% -- FLEX: [CORE]
}

\centering
\Large\bfseries
Can you make the model 4$\times$ smaller\\[0.2cm]
and get a 4$\times$ speedup?
\end{frame}

% --- Slide 3.2: Wall 13 --- Compression ---
\begin{frame}{Wall 13: The Fidelity Wall}
\note{[2 min] ``Storage always shrinks. But inference speedup
depends on the method. This distinction trips up everyone.''
% -- FLEX: [CORE]
}

\small
\wallbox{The Fidelity Wall}{
\[
  \text{Compression}_{\text{quant}} = \frac{32}{\text{bits}},
  \qquad
  \text{Compression}_{\text{prune}} = \frac{1}{1 - \text{sparsity}}
\]
}

\vspace{0.3cm}
\begin{tabular}{lccc}
\toprule
\textbf{Method} & \textbf{Storage} & \textbf{Speedup} & \textbf{Accuracy} \\
\midrule
FP32 $\to$ FP16 & 2$\times$ & 2$\times$ & $\sim$0\% loss \\
FP16 $\to$ INT8  & 2$\times$ & 1.5--2$\times$ & $<$1\% loss \\
FP16 $\to$ INT4  & 4$\times$ & 2--3$\times$ & 2--5\% loss \\
50\% unstructured prune & 2$\times$ & \alert{1$\times$ (no speedup!)} & 1--3\% loss \\
50\% structured prune & 2$\times$ & $\sim$2$\times$ & 2--5\% loss \\
2:4 N:M sparsity & 2$\times$ & 2$\times$ & 1--2\% loss \\
\bottomrule
\end{tabular}

\vspace{0.2cm}
\alert{Unstructured pruning saves storage but gives zero GPU speedup.\\
Only structured patterns accelerate hardware execution.}
\end{frame}

% --- Slide 3.3: Quantization Deep Dive ---
\begin{frame}{Quantization: Trading Bits for Speed}
\note{[2 min] Walk through the precision ladder.
``If you quantize Llama-3 8B from FP16 to INT4, how much memory?''
Expected: 8B * 0.5 bytes = 4 GB. Down from 16 GB.
% -- FLEX: [CORE]
}

\small
\begin{columns}[T]
\begin{column}{0.55\textwidth}
\textbf{The precision ladder:}

\vspace{0.2cm}
\begin{tikzpicture}[scale=0.8, >=Stealth]
  \foreach \i/\label/\bits/\color in {
    0/FP32/32/errorfill,
    1/FP16\slash BF16/16/routingorange,
    2/INT8\slash FP8/8/computeblue,
    3/INT4/4/datagreen} {
    \fill[\color, draw=midgray] (0, 3-\i) rectangle ({0.15*\bits}, 3.6-\i);
    \node[right, font=\footnotesize] at ({0.15*\bits + 0.2}, 3.3-\i) {\label};
  }
  \draw[->, thick] (0, -0.5) -- (5.5, -0.5) node[right, font=\scriptsize] {Size};
  \node[font=\scriptsize, midgray] at (2.5, -1) {$\longleftarrow$ smaller is better};
\end{tikzpicture}
\end{column}
\begin{column}{0.42\textwidth}
\textbf{Llama-3 8B memory:}
\scriptsize
\begin{tabular}{lc}
\toprule
Precision & Weight Size \\
\midrule
FP32 & 32 GB \\
FP16 & 16 GB \\
INT8 & 8 GB \\
INT4 & 4 GB \\
\bottomrule
\end{tabular}

\vspace{0.3cm}
\normalsize
At INT4, Llama-3 8B fits in\\
a \textbf{laptop GPU} (6 GB).
\end{column}
\end{columns}
\end{frame}

% --- Slide 3.4: Live Demo --- CompressionModel ---
\begin{frame}[fragile]{Live Demo: Quantization Impact}
\note{[2 min] Run CompressionModel live. Show storage savings
and speedup side by side.
% -- FLEX: [CORE]
}

\small
\begin{lstlisting}
from mlsysim.core.solver import CompressionModel
comp = CompressionModel()
for bits in [16, 8, 4]:
    r = comp.solve(llama, hw, method="quantization",
                   target_bitwidth=bits)
    print(f"INT{bits}: size={r.compressed_size:~P}  "
          f"speedup={r.inference_speedup:.1f}x")
\end{lstlisting}

\vspace{0.2cm}
\begin{exampleblock}{What to observe}
  \begin{itemize}
    \item Storage shrinks linearly with bit reduction
    \item Speedup follows storage for quantization (structured by nature)
    \item Accuracy degrades modestly at INT8, more at INT4
  \end{itemize}
\end{exampleblock}
\end{frame}

% --- Slide 3.5: Structured vs Unstructured Pruning ---
\begin{frame}[fragile]{Pruning: The Structure Matters}
\note{[2 min] ``Unstructured pruning zeros out individual weights.
The matrix is still the same shape, so the GPU does the same
number of operations. No speedup! Structured pruning removes
entire rows/columns, physically shrinking the matrix.''
WARN: Students assume any compression = speedup.
% -- FLEX: [CORE]
}

\small
\begin{columns}[T]
\begin{column}{0.48\textwidth}
\textbf{Unstructured}
\begin{itemize}\setlength\itemsep{2pt}
  \item Zero out individual weights
  \item Matrix shape unchanged
  \item GPU does same work (skip zeros? nope)
  \item \alert{Storage savings only}
\end{itemize}
\end{column}
\begin{column}{0.48\textwidth}
\textbf{Structured / N:M}
\begin{itemize}\setlength\itemsep{2pt}
  \item Remove entire rows/columns, or\\
        2:4 pattern (Ampere+)
  \item Physically smaller matrices
  \item GPU hardware support (2:4 $\to$ 2$\times$)
  \item \textbf{Real speedup}
\end{itemize}
\end{column}
\end{columns}

\vspace{0.3cm}
\begin{lstlisting}
for stype in ["unstructured", "structured", "n_m"]:
    r = comp.solve(llama, hw, method="pruning",
            sparsity=0.5, sparsity_type=stype)
    print(f"{stype:>14}: {r.inference_speedup:.1f}x")
\end{lstlisting}
\end{frame}

% --- Slide 3.5b: Predict --- What Does INT4 Change? ---
\begin{frame}{Predict: What Does INT4 Change?}
\note{[1 min] Quick poll. Most say "latency gets better." The real answer is fleet size.
Turn to your neighbor: did you get the same answer? Why or why not? 60 seconds.}
\centering
\Large
\textbf{You quantize Llama-3 70B from FP16 to INT4.}\\[0.5cm]
\normalsize
What is the BIGGEST impact on your serving infrastructure?\\[0.3cm]
\begin{enumerate}[(A)]
  \item Inference latency drops by 4$\times$
  \item Model quality degrades significantly
  \item \textbf{You need half as many GPUs}
  \item Memory bandwidth becomes the bottleneck
\end{enumerate}
\end{frame}

% --- Slide 3.6: Predict Before You Peek #4 ---
\begin{frame}{Predict: INT4 Llama-3 8B on H100}
\note{[2 min] Predict-before-reveal. If you quantize to INT4 at bs=1,
is it still memory-bound? Yes! Load time = 4/3350 = 1.2 ms.
Compute still ~0.03 ms. Still memory-bound but 4x faster.
Turn to your neighbor: did you get the same answer? Why or why not? 60 seconds.
% -- FLEX: [CORE]
}

\centering
\Large\bfseries
Quantize Llama-3 8B to INT4.\\
Run inference at batch size 1 on H100.\\[0.5cm]
\normalsize
Is it still memory-bound?\\[0.3cm]

\pause

\small
\begin{tabular}{lcc}
\toprule
              & \textbf{FP16} & \textbf{INT4} \\
\midrule
Weight size   & 16 GB & 4 GB \\
Load time     & 4.8 ms & 1.2 ms \\
Compute time  & 0.03 ms & 0.03 ms \\
\midrule
Bottleneck    & Memory & \textbf{Still Memory!} \\
Decode speedup & --- & \textbf{4$\times$} \\
\bottomrule
\end{tabular}

\vspace{0.3cm}
\pause
\alert{INT4 gives 4$\times$ faster decode, but the GPU is still memory-bound.\\
The memory wall is that deep.}
\end{frame}

% --- Slide 3.7: Exercise 3 ---
\begin{frame}[fragile]{Exercise 3: Compression Tradeoffs}
\note{[5 min] Compare INT8 quantization vs 50\% structured pruning
for Llama-3 8B on H100. Which gives better speedup per accuracy loss?
Expected: INT8 wins (~2x speedup, <1\% loss vs 2-5\% loss).
Turn to your neighbor: did you get the same answer? Why or why not? 60 seconds.
% -- FLEX: [CORE]
}

\centering
\Large\bfseries
Hands-On Exercise\\[0.5cm]
\normalsize

\textbf{Question:} For Llama-3 8B on H100, which is the better deal?
\begin{enumerate}
  \item INT8 quantization
  \item 50\% structured pruning
\end{enumerate}
Compare speedup per accuracy point lost.

\vspace{0.3cm}
\begin{lstlisting}
r1 = comp.solve(llama, hw, method="quantization",
                target_bitwidth=8)
r2 = comp.solve(llama, hw, method="pruning",
        sparsity=0.5, sparsity_type="structured")
print(f"Quant: {r1.inference_speedup:.1f}x / "
      f"{abs(r1.accuracy_delta):.1%} loss")
print(f"Prune: {r2.inference_speedup:.1f}x / "
      f"{abs(r2.accuracy_delta):.1%} loss")
\end{lstlisting}
\end{frame}

% --- Slide 3.7b: Compression Changes Fleet Architecture ---
\begin{frame}{Compression Changes Fleet Architecture}
\note{[3 min] This is the ``aha'' that compression is architecture, not optimization.
The punchline: INT4 halves your GPU count AND your electricity bill.}

\small
\textbf{Llama-3 70B Serving Fleet:}

\vspace{0.3cm}
\begin{tabular}{@{}lrrr@{}}
\toprule
Precision & Model Size & GPUs Needed & Annual Cost \\
\midrule
FP16 & 140 GB & 4 (TP=4) & \$480K \\
INT8 & 70 GB & 2 (TP=2) & \$240K \\
INT4 & 35 GB & 1 & \$120K \\
\bottomrule
\end{tabular}

\vspace{0.3cm}
\textbf{INT4 doesn't just improve latency --- it eliminates 3 GPUs per replica.}\\
At 100 replicas for 1000 QPS: that's \textbf{300 fewer GPUs} and \textbf{\$36M saved per year}.

\vfill
\centering
\small\textcolor{gray}{This is why quantization is a Day 1 architectural decision, not a Day 100 optimization.}
\end{frame}

% --- Slide 3.8: Part 3 Key Takeaway ---
\begin{frame}{Part 3: Key Takeaway}
\note{[1 min] One sentence. Repeat.
% -- FLEX: [CORE]
}

\centering\Large

\textbf{Storage savings $\neq$ inference speedup.}\\[0.5cm]

\normalsize
\begin{itemize}
  \item Quantization gives both storage and speed gains.
  \item Unstructured pruning gives storage only --- zero GPU speedup.
  \item N:M sparsity (2:4) is the hardware-friendly middle ground.
  \item Even at INT4, LLM decode is \textit{still} memory-bound.
  \item \texttt{CompressionModel.solve()} quantifies the full tradeoff.
\end{itemize}
\end{frame}

% --- Roadmap: After Lunch ---
\begin{frame}{Roadmap: Afternoon Session}
\note{[1 min] Re-energize the room. ``Welcome back. The morning was about
single-node physics. The afternoon is about fleets, money, and carbon.''}

\centering\small
\begin{tabular}{rll}
\toprule
\textbf{Time} & \textbf{Part} & \textbf{Status} \\
\midrule
9:00--12:00  & Parts 0--3: Single Node   & \checkmark Done \\
\midrule
\rowcolor{crimson!12}
1:00--2:15   & \textbf{Part 4: Going Distributed} & \textbf{$\leftarrow$ You are here} \\
2:30--3:15   & Part 5: Economics \& Sustainability & \\
3:15--3:45   & Part 6: Design Space Exploration & \\
3:45--4:15   & Part 7: TinyML to Frontier & \\
4:15--4:45   & Part 8: Advanced Topics & \\
4:45--5:00   & Part 9: Wrap-Up \& Capstone & \\
\bottomrule
\end{tabular}
\end{frame}

% =============================================================================
% PART 4: GOING DISTRIBUTED (15 slides)
% =============================================================================


\end{document}
